\scanpage{136}
Le résultat précédent nous permet de démontrer le résultat auxiliaire suivant, qui nous permettra de démontrer la formule (6) en \deleted{\unsure{toute}} généralité :

\uline{Lemme} Soient $A$ et $B$ deux opérateurs $\in\mathcal L^{2m}$ \unsure{($\mathcal H_2$)|($\mathcal H$)|($\mathfrak H_2$)|($\mathfrak H$)|($H_2$)|($H$)},
\margin{posons\;\rule[-2.1em]{0.4pt}{3.3em}}
$H=A^*A,\ K=B^*B,\ L=(AB)^*(AB)$.
Soient \deleted{\unsure{$L$ soit|et soit}} $m$ et $n$ des entiers $\geq1$.
Alors \deleted{\unsure{$L^m$ est, $H$|$L^n,H^n$|$L^m\in\mathcal L^1,H$}} \inserted{$L^m,H^m$} et $K^m$ sont opérateurs à trace, et on a
\[
\alpha_n\bigl(|\deleted{\unsure{A}}L|^m\bigr)
\leq\alpha_n\bigl(\deleted{\unsure{H}}H^m\bigr)\,\alpha_n\bigl(\unsure{\deleted{\unsure{K}}K^m|\lvert K^m}\bigr)
\tag*{(12)}
\]
Le fait que $L^m,H^m,K^m$ soient opérateurs à trace résulte \deleted{\unsure{enfin d'un des précédents}} \inserted{aussitôt de ce qui précède}. (12) s'écrit alors, \deleted{\unsure{cela|et|entiers}}
\deleted{\unsure{par contre s'écrit en effet en posant}}
\inserted{en posant} $\mathscr A=\underline{\wedge} A,\ \mathscr B=\underline{\wedge} B$
\inserted{$\mathscr H=\mathscr A^*\mathscr A,\ \mathscr K=\mathscr B^*\mathscr B,\ \mathscr L=(\mathscr A\mathscr B)^*(\mathscr A\mathscr B)$}
\[
\operatorname{Tr}\,|\deleted{\unsure{\mathscr A\mathscr B}}\mathscr L|^{\unsure{m|2m}}\leq
\operatorname{Tr}\,\deleted{\unsure{\mathscr H}}\mathscr H^m\ \operatorname{Tr}\,\deleted{\unsure{\mathscr B}}\unsure{\lvert\mathscr K\rvert^m|\mathscr K^m},
\]
\deleted{\illegible{2}}
\[
\deleted{\unsure{\bigl(\|\mathscr A\mathscr B\|_m\bigr)^m\leq
\bigl(\|\mathscr A\|_m\bigr)^m\bigl(\|\mathscr B\|_m\bigr)^m|\bigl(\|L\|_m\bigr)^m\leq\bigl(\|H\|_m\bigr)^m\bigl(\|K\|_m\bigr)^m}}\quad
\deleted{\text{\unsure{ce qui}}}
\]
\deleted{$\|\mathscr A\mathscr B\|_m\leq\|\mathscr A\|_m\|\mathscr B\|_m$. Or, \unsure{on a vu qu'on a même}
$\|AB\|_m\unsure{<|\leq}\|A\|_{\frac m2}\|B\|_{\frac m2}$}

ce qui nous ramène au cas où $n=1$. On doit alors démontrer
\[
\unsure{\operatorname{Tr}\||\operatorname{Tr}\lvert|\operatorname{Tr}[}\,(ABB^*A^*)\mathinner{\cdotp\cdotp\cdotp-}(ABB^*A^*)\|_2
\leq\bigl(\|\deleted{\unsure{H}}A^*A\|_m\bigr)^m\bigl(\|B^*B\|_m\bigr)^m
\]
Or, le premier membre \deleted{\unsure{est la trace-norme}} \inserted{est aussi égal à la}
\deleted{\unsure{d'un produit de $4m$ facteurs, est en effet}} trace du produit de $2m$ facteurs
\deleted{\unsure{majoré par}}
\[
(A^*A)(BB^*)\ldots(BB^*),
\]
donc au plus égal \unsure{à|:} la trace-norme de ce produit, donc au plus égal \unsure{à|:}
$\bigl(\|A^*A\|_{2m}\bigr)^m\bigl(\|BB^*\|_{2m}\bigr)^m$
d'après la formule $\|A_1\ldots A_{2m}\|_1\leq\|A_1\|_{2m}\ldots\|A_{2m}\|_{\unsure{2m}}$

Or, on a $\|A^*A\|_{2m}\leq\|A^*A\|_m$, et
$\|BB^*\|_{2m}\leq\|B^*B\|_{2m}\leq\|B^*B\|_m$, d'où résulte le lemme.
